\paragraph{Choice of the light generators.}
On each run use the scalar \(s_W\ne0\) assigned to its
distinguished prime and light weight, and put
\(M_0=g_0g_1\). If \(D_F\) denotes the polynomial degree
of the light weight, then \(D_F\le\operatorname{DF}(\kappa)\)
and \(2\operatorname{DF}(\kappa)+2<g_i\) for \(i=0,1\).

In particular 
\[
  M_0\ge4D_F+4.
\]
 The half-level theta
module at this weight is generated by harmonic polynomial
differential rules. Its short generating family includes
the translates of a single block determinant rule of
degree \(D_F\). Multiply the entire family by \(s_W\).
The scaled theta lattice is inside the integral cusp
lattice by the weight and level integrality comparison.
The count and upper norms of its generators are those of
the unscaled theta family with the norm of \(s_W\) included.
In the tables this gives
\(\log N\le n_B\), 
\[
  \log N\le P_{rw}(S_L-1),
\]
 and the
mean log bound \(G_c S_L\) for the maximum of their
norms over the generator list. In particular the number
of terms in the sum can be charged once using \(\log N\).
\label{cor:slice-light-1}

\begin{auxlemma}[Interpolation on the quotient by the short periods]
\label{lem:light-thick-interpolation}
Let \(Y=\mathbb C^h/\Gamma\) have a principal Riemann form
\(\operatorname{Im}H\) and its Gaussian theta metric. Use Haar
probability in the squared section norm. Let \(M\ge1\), \(c\ge0\),
and suppose
\[
 \sqrt{\pi M}\min_{0\ne\gamma\in\Gamma}\|\gamma\|\ge R_0,\qquad
 R_0^2\ge\max(270,90(h+c)).
\]
For \(h=0\) the lattice condition is empty. For each unit direction
\(y\), or for the constant polynomial when \(c=0\), there is a
level-\(M\) theta section \(\beta\) whose Taylor polynomial through
degree \(c\) is \(p(z)=H(z,y)^c\). It satisfies
\[
 \|\beta\|_2^2\le2M^{-h}(\pi M)^{-c}c!,\qquad
 \sup_z |\beta(z)|^2e^{-\pi M\|z\|^2}\le4(\pi M)^{-c}c!.
\]
The constant jet \(p=1\) can likewise be prescribed through any
order \(c\) in this range. Its bounds are
\(\|\beta\|_2^2\le2M^{-h}\) and Gaussian supremum squared at most four.
\end{auxlemma}
\begin{proof}
Put the Gaussian norm
\(\|Q\|_F^2=\int_{\mathbb C^h}|Q(z)|^2e^{-\pi M\|z\|^2}dz\)
on polynomials of degree at most \(c\). Let \(U_\gamma Q\) be
\(Q(z-\gamma)\) multiplied by the factor of automorphy
\(\epsilon_\gamma\exp(\pi M H(z,\gamma)-\pi M\|\gamma\|^2/2)\),
with the semicharacter of the line. This operator preserves the
Gaussian norm. Define \(SQ=\sum_{\gamma\in\Gamma}U_\gamma Q\).
Polynomial growth times Gaussian decay gives normal convergence on
compact sets, also after each fixed differentiation. The cocycle
identity makes \(SQ\) a section.

Let \(TQ\) be the Taylor polynomial of \(SQ\) through degree \(c\).
Rotation in each complex coordinate makes distinct monomials
orthogonal in the Gaussian integral. For any polynomial \(Q'\) of
degree at most \(c\) it follows that
\[
 \langle (T-1)Q,Q'\rangle_F
     =\sum_{\gamma\ne0}\langle U_\gamma Q,Q'\rangle_F .
\]
One can first integrate on balls and then increase the radius.
Normal convergence and Gaussian integrability justify the limit.
Write \(t=\|\gamma\|/2\). The two sets \(\|z\|\ge t\) and
\(\|z-\gamma\|\ge t\) cover the integration domain. For any
polynomial of degree at most \(c\), dilation and orthogonality of
its homogeneous terms give
\[
 \int_{\|z\|\ge t}|Q(z)|^2e^{-\pi M\|z\|^2}dz
 \le 2^{h+c}e^{-\pi M t^2/2}\|Q\|_F^2.
\]
Cauchy--Schwarz on the two sets therefore bounds the absolute
value of each inner product by
\(2^{1+(h+c)/2}e^{-\pi M\|\gamma\|^2/16}\|Q\|_F\|Q'\|_F\).

For the scaled lattice \(\sqrt{\pi M}\,\Gamma\), disjoint balls
of radius \(R_0/2\) give at most \((3u)^{2h}\) nonzero points
of norm at most \(R_0u\) for \(u\ge1\). Put \(a=R_0^2/16\).
Integrate this counting bound against \(2au e^{-au^2}\,du\).
Since \(2\log u\le u^2-1\) and \(a\ge2h\), this gives
\[
 \sum_{\gamma\ne0}e^{-\pi M\|\gamma\|^2/16}
 \le 3^{2h}\frac{a}{a-h}e^{-a}\le2\,9^he^{-a}.
\]

The product of this bound and \(2^{1+(h+c)/2}\) is less than
\(1/2\), by the two lower bounds for \(R_0^2\). In particular
\(T\) is invertible and \(\|T^{-1}\|\le2\). Set \(Q=T^{-1}p\),
\(\beta=SQ\). Unfolding the section norm against one copy of
the sum gives
\[
 \|\beta\|_2^2=\langle SQ,Q\rangle_F
        =\langle p,Q\rangle_F\le2\|p\|_F^2 .
\]

Here the Euclidean covolume is one because the Riemann form is
principal in the compatible unitary metric. A unitary change of
variables and a radial Gaussian integral give
\(\|p\|_F^2=M^{-h}(\pi M)^{-c}c!\).

For the point estimate, the evaluation kernel on the covering
space is \(M^h\exp(\pi M H(z,w))\). Gaussian integration of the
Taylor series after translation to \(w\) gives its reproducing
identity. Its periodization gives the evaluation kernel on \(Y\)
by unfolding. On the diagonal, in the Gaussian metric, its modulus
is at most \(M^h\sum_\gamma e^{-\pi M\|\gamma\|^2/2}\le2M^h\)
by the preceding counting bound. The reproducing identity and
Cauchy--Schwarz give the asserted supremum bound.
\end{proof}

\begin{auxlemma}[Lifting a jet in a direction orthogonal to the short periods]
\label{lem:light-thick-lift}
Use the parameters of Lemma~\ref{lem:joint-theta-parameters}.
Let \(I\) be a saturated isotropic period lattice on which the
semicharacter is trivial. Write \(W=I\otimes\mathbb R\),
\(V_B=(W+iW)^\perp\), and \(U\subset W\) for the projections onto
\(W\) of the imaginary coordinates of all periods along \(W+iW\).
Assume
\[
 \operatorname{rank}I\le d=2fm,\qquad
 \|u-u'\|\ge\sqrt B\quad(u\ne u',\ u,u'\in U),
\]
and that the induced
principally polarized quotient torus in \(V_B\) has scaled first
minimum at least \(R_0\) as in Lemma~\ref{lem:light-thick-interpolation}.
For each unit \(y\in V_B\) in one embedding space and
\(1\le c\le D\), there is a level-\(M\) section \(s\) on the original
torus with
\[
 \begin{gathered}
 \left|d^cs(0)[v_1,\ldots,v_c]-c!\prod_lH(v_l,y)\right|
         \le \varepsilon c!,\qquad \|s\|_2^2\le K_6(c!)^2,\\
 \varepsilon=7\,2^g(4.3)^d e^{-T_0},
 \qquad K_6=6(2d)^{d/2}
 \end{gathered}
\]
for unit directions \(v_l\) in this embedding. An unused column admits
a section with \(|s(0)|\ge1/2\) and \(\|s\|_2^2\le K_6\).
\end{auxlemma}
\begin{proof}
Take orthonormal real coordinates on \(W\), write \(z\) for their
complexification and \(w\) for the \(V_B\)-coordinate. Lift the
section \(\beta\) of Lemma~\ref{lem:light-thick-interpolation} by
\[
 F_0(z,w)=\exp\!\left(\frac{\pi M}{2}\sum_j z_j^2\right)\beta(w).
\]

The multiplier under a period with imaginary coordinate zero is
precisely the multiplier of the original line. In fact the real
translation on \(W\) contributes the difference of the two
quadratic exponents, and the remaining multiplier is the one
defining the quotient line. The trivial semicharacter on \(I\)
makes the quotient multiplier independent of its lift.

Sum the unitary translates of \(F_0\) over the remaining period
classes, indexed by \(u\in U\), to form \(s\).
A translated term, in the Gaussian metric, has absolute value
\[
 e^{-\pi M\|\operatorname{Im}z-u\|^2}
       |\beta(w-w_u)|e^{-\pi M\|w-w_u\|^2/2},
\]
where \(w_u\) is the perpendicular component of the representative.
The Gaussian and the lattice separation imply normal convergence
on compact sets; reindexing gives the full theta transformation law.

Integration over \(W/I\) kills cross terms because their characters
have different frequencies \(M u\). Average next over the quotient
torus in \(w\). Every translate of \(\beta\) has the same squared
norm. The remaining Gaussian sum is bounded uniformly in
\(t=\operatorname{Im} z\): choose a nearest \(u_0\in U\), so
\(\|t-u\|\ge\|u-u_0\|/2\), and apply the preceding packing integral
in real dimension at most \(d\). With \(T_0=MB\) this gives
\[
 \sum_{u\in U}e^{-2\pi M\|t-u\|^2}
 \le1+2\,3^d e^{-\pi T_0/2}\le 2,\qquad
 \|s\|_2^2\le 2\|\beta\|_2^2 .
\]

The last comparison uses Haar probability on every quotient.
In particular \(\|s\|_2^2\le4M^{-\dim_{\mathbb C}V_B}
(\pi M)^{-c}c!\le K_6(c!)^2\).

It remains to bound the derivatives of the terms \(u\ne0\).
In the holomorphic frame the Gaussian bound and
\(\|u-t\|^2\ge\|u\|^2/2-\|t\|^2\) yield
\[
 |s(v)-F_0(v)|\le
 2\sqrt{c!}(\pi M)^{-c/2}
 e^{3\pi M\|v\|^2/2}\,2\,3^d e^{-\pi T_0/2}.
\]

To differentiate in the \(c\) prescribed unit directions, restrict
to \(v=\sum_{l=1}^c z_lv_l\), with all \(|z_l|\le(M(c+1))^{-1/2}\).
Here \(M\|v\|^2\le c\), and the coefficient of
\(z_1\cdots z_c\) is the desired derivative. Cauchy's inequality,
followed by \(c!\ge(c/e)^c\), bounds the error divided by \(c!\) by
\[
 4\,3^d e^{-\pi T_0/2}
 e^{3\pi c/2}\frac{((c+1)/\pi)^{c/2}}{\sqrt{c!}}
 \le 4\,3^d e^{6c-\pi T_0/2}\le\varepsilon .
\]

Indeed \(T_0\ge28D\) by \eqref{eq:joint-theta-log-cost}.
The zero mode has the jet of \(H(w,y)^c\), since its lower
Taylor coefficients vanish and the exponential factor is one
at zero. This proves the derivative comparison. With \(c=0\)
the same construction starting from the constant jet gives
\(|s(0)-1|\le4\,3^de^{-\pi T_0/2}\le1/2\) and the asserted norm.
\end{proof}

\begin{auxlemma}[Separation of harmonic jets at the marked level]
\label{lem:ideal-lattice-jet}
Let \(\Lambda\) be the marked self-dual Hermitian lattice at a complex
frame and \(\mathfrak m\bar{\mathfrak m}=(M_0)\) its chosen half-level
ideal. Let \(D\) be a nonzero constant-coefficient differential operator
of order \(d\) in the \(n\) columns of the frame, with \(M_0\ge4d+4\).
Write \(\theta_\nu\) for the theta sections with the characters of the
marked quotient in those columns. In the theta trivialization at the
origin,
\[
             (D\theta_\nu(0))_\nu\ne0 .
\]
\end{auxlemma}
\begin{proof}
Let \(A'\to A\) be the isogeny of complex abelian varieties given by
the inclusion of period lattices \(\mathfrak m\Lambda\subset\Lambda\).
Take its product \(\pi:(A')^n\to A^n\) and denote by \(L\) the product
of the principal theta lines on \(A^n\). If \(B\) is the alternating
Riemann form on the original lattice, the Riemann form on the ideal
lattice is \(B/M_0\) for its principal polarization. Indeed the
\(B\)-dual of \(\mathfrak m\Lambda\) is
\(\bar{\mathfrak m}^{-1}\Lambda\), so its dual for \(B/M_0\) is
\[
        M_0\bar{\mathfrak m}^{-1}\Lambda=\mathfrak m\Lambda .
\]

It follows that \(\pi^*L\) is an \(M_0\)-th power of an ample line
up to a degree-zero twist. Absorb the twist in one ample factor.
Since \(M_0\ge d+2\), grouping factors gives a product of \(d+2\)
ample line bundles. Theorem~1 of Bauer and Szemberg gives the
surjection \cite{bauer-szemberg-jets}
\[
 H^0((A')^n,\pi^*L)\longrightarrow
 (\pi^*L)_0/\mathfrak m_0^{d+1}(\pi^*L)_0 .
\]

Here \(\mathfrak m_0\) denotes the maximal ideal at the origin.
The functional induced by \(D\) on this jet space is nonzero. In local
coordinates a monomial of degree \(d\) paired with a nonzero
top coefficient has nonzero image in characteristic zero.

Write \(H=\ker\pi\). Character decomposition of the finite cover
over \(\mathbb C\) and the projection formula give
\[
 H^0((A')^n,\pi^*L)
  =\bigoplus_{\nu\in\widehat H} H^0(A^n,L\otimes P_\nu),
\]
where \(P_\nu\) is the degree-zero line of character \(\nu\).
Each summand is one-dimensional. The polarization is principal,
so \(L\otimes P_\nu\) is a translate of \(L\). The half-level marking
indexes these translates by the characters of \(H\). Their pullbacks
in the common theta trivialization are precisely the \(\theta_\nu\),
up to nonzero constant factors. Thus \(D\) cannot vanish on all the
summands. Changing an orthonormal basis makes an invertible linear
change of variables on the jet and preserves this conclusion.
\end{proof}

\begin{auxlemma}[Selecting the column directions of a block]
\label{lem:light-block-directions}
In an embedding space of dimension \(N\le r\), take an orthonormal
basis \(w_1,\ldots,w_N\) and \(i\le N\) orthonormal row
directions \(e_a\). Suppose that for \(1\le j\le t\) the available
direction \(y_j\) satisfies
\[
 H(v,y_j)=H(v,w_j)+\sum_{l\le t}E_{jl}H(v,w_l),
 \qquad \sum_{l\le t}|E_{jl}|\le\frac1{30t^{3/2}},
\]
for all row directions \(v\). For \(j>t\) take \(y_j=w_j\).
One can choose \(i\) distinct indices \(j_b\) so that
\[
 \left|\det(H(e_a,y_{j_b}))_{a,b}\right|\ge 2^{-r/2}/2 .
\]
\end{auxlemma}
\begin{proof}
The \(i\)-by-\(N\) matrix \(C_{aj}=H(e_a,w_j)\) satisfies
\(CC^*=1\). Cauchy--Binet shows that the squared absolute values
of its \(i\)-rowed minors sum to one. The maximum modulus \(\Delta\)
is thus at least \(\binom Ni^{-1/2}\ge2^{-r/2}\). Choose its columns.
At most \(t\) of them change. Expand the determinant multilinearly
in these columns. Each minor of the original matrix appearing in
an error term has modulus at most \(\Delta\); repeated columns
give zero. For \(t>0\) the sum of the absolute errors is therefore
bounded by
\[
 \Delta\bigl((1+(30t^{3/2})^{-1})^t-1\bigr)
 \le\Delta(e^{1/30}-1)<\Delta/2 .
\]
For \(t=0\) there is no error.
\end{proof}

\begin{auxlemma}[The short periods and their dual lattice]
\label{lem:light-isotropic-periods}
Use the principally polarized CM torus \(V/L\) of
Lemma~\ref{lem:eta-period-decay} and the parameters of
Lemma~\ref{lem:joint-theta-parameters}. Write \(A=\operatorname{Im}H\)
and \(\lambda=\lambda_1(L)\). There is a saturated
\(\mathcal O_E\)-submodule \(I\subset L\) with isotropic real
span \(W\). Write \(U(x)\in W\) for the vector representing
\(w\mapsto A(x,w)\) in the Euclidean metric and \(X=U(L)\).
Then
\[
 I=W\cap L,\qquad X=I^*,\qquad \dim_{\mathbb R}W\le d=2fm,
 \qquad I=0\ \Longleftrightarrow\ \lambda\ge h .
\]
The induced lattice on the complex orthogonal complement of
\(W+iW\) has minimum at least \(h\). If \(I\ne0\), then
\begin{equation}\label{eq:light-dual-successive}
 \lambda_1(X)^2\ge B,\qquad
 \lambda_{\dim W}(X)\le d^2/\lambda .
\end{equation}
\end{auxlemma}
\begin{proof}
Start with \(W=0\), and keep \(W\) isotropic, rational with
respect to \(L\), and stable under \(E\otimes\mathbb R\).
Maintain the inequality \(\|U(y)\|>(\Omega h)^{-1}\) whenever
\(y\in L\) has nonzero restriction to \(W\). If some
\(x\in L\setminus W\) has distance less than \(h\) from \(W\),
write \(x=w+z\) with \(w\in W\), \(z\perp W\). Then
\[
 \|U(x)\|^2=A(z,U(x))\le\|z\|\,\|U(x)\|,
 \qquad \|z\|<h\le(\Omega h)^{-1}.
\]

Thus \(A(x,W)=0\). Take integral \(\omega_j\) spanning \(E\)
over \(\mathbb Q\) with embedding norms bounded by \(\Omega\), as in
Lemma~\ref{lem:eta-period-decay}. Integrality and the Rosati
identity give
\[
 |A(x,\omega_jx)|=|A(z,\omega_jz)|\le\Omega\|z\|^2<1.
\]

Hence \(A(x,Ex)=0\). Enlarge \(W\) by the real span of
\(\mathcal O_E x\) and intersect with \(L\). This preserves
isotropy. For a nonzero restriction of \(A(y,\cdot)\) to the
larger space, its norm dominates the previous norm if the
restriction was already nonzero. Otherwise it has nonzero integral
value at some \(\omega_jx\). Its representing vector \(u\) on the
new space is orthogonal to the old one, so
\[
       1\le |(u,\omega_j z)|<\|u\|\Omega h .
\]

This proves the maintained inequality. The dimension increases at
each enlargement. At termination each lattice point outside \(W\)
has distance at least \(h\) from \(W\). In particular for
\(y\in L\) with \(A(y,W)=0\), the component perpendicular
to \(W\) lies in \((W+iW)^\perp\) and has that length unless
\(y\in I\). These components form the quotient lattice. Its alternating
form is again unimodular: extend a basis of the primitive
isotropic sublattice \(I\) to an integral symplectic basis of \(L\).
This can be done one primitive vector at a time, choosing a dual
vector of pairing one and splitting off their hyperbolic plane.
The remaining symplectic pairs give a basis of \(I^{\perp_A}/I\).
In the quotient Gaussian coordinates its scaled minimum satisfies
\[
     \sqrt{\pi M}\,h=R_0\sqrt{\pi f\rho^{21}}\ge R_0 .
\]


The initial enlargement occurs exactly when \(\lambda<h\).
The rational span of \(I\) has an \(E\)-dimension \(a\).
Isotropy implies \(W\cap iW=0\), so \(W+iW\) realizes its
complexification inside \(V\). Each embedding occurs there with
dimension \(a\). Since some embedding block of \(V\) has
dimension at most \(m\), we obtain \(\dim_{\mathbb R}W=2fa\le d\).
The theta semicharacter is trivial on \(I\), since
\(A(k,uk)=0\) there for the half-trace element \(u\).
Every integral functional on \(I=L\cap W\) extends to \(L\).
Perfectness of \(A\) gives \(U(L)=I^*\).
Finally
\[
  B\Omega^2h^2=4f^4/\rho^5\le1,\qquad
  \lambda_{\dim W}(I^*)\lambda_1(I)\le\dim W\le d
\]
by Euclidean transference. With \(\lambda_1(I)\ge\lambda\) this
implies \eqref{eq:light-dual-successive}. \cite{banaszczyk}
\end{proof}

\begin{auxlemma}[Frequencies in one CM embedding]
\label{lem:light-triangular-frequencies}
In Lemma~\ref{lem:light-isotropic-periods}, let \(V_\sigma\) be
the space for one embedding and \(A_\sigma=(W+iW)\cap V_\sigma\)
have dimension \(a>0\). There exist an orthonormal basis
\(e_1,\ldots,e_a\) of \(A_\sigma\), frequencies \(u_s\in X\),
and real \(R_s>0\) such that, with \(p_\sigma\) the orthogonal
projection to \(V_\sigma\),
\begin{align}
 p_\sigma u_s/R_s &=e_s+\sum_{i<s}t_{si}e_i,
 &\sum_{i<s}|t_{si}|^2&\le (30a^2)^{-2},\label{eq:light-triangular-projection}\\
 1&\le2\pi R_s, & \|u_s\|&\le 3R_s/2,\nonumber\\
 |2(u_s,y)|&\le 90m^2d^2\rho\,\|y\|^2\quad(y\in X),
 &\|u_s\|^2&\le 8100m^4d^4\rho^7 P.\label{eq:light-frequency-cost}
\end{align}
\end{auxlemma}
\begin{proof}
Put \(K=d^2\rho\). Let \(A'\) be the span of the projections
already selected. Choose a shortest \(v\in X\) with
\(p_\sigma v\notin A'\), and put \(\mu=\|v\|\). The choice
exists by discreteness and \eqref{eq:light-dual-successive}, which
also gives
\[
            1\le\mu\le d^2/\lambda .
\]
Let \(Z\subset W\) be the real span of the \(\mathcal O_E\)-multiples
of frequencies of length at most \(\mu\). It is spanned by the
products of these frequencies with the \(\omega_j\), hence by
vectors of \(X\) of length at most \(\Omega\mu\). Here stability
of the dual follows from the Rosati identity. The spaces \(Z\)
are nested, since the successive shortest lengths cannot decrease.
The projections chosen at earlier steps lie in \(p_\sigma Z\).

On the \(\{\sigma,\bar\sigma\}\)-summand of the real CM
decomposition of \(Z\), the real CM algebra contains \(J=(i,-i)\).
Isotropy, applied to \(z,Jz\), gives
\[
 0=A(z,Jz)=\pm\big(\|p_\sigma z\|^2-\|p_{\bar\sigma}z\|^2\big),
 \qquad \|p_\sigma z\|^2=\tfrac12\|z\|^2 .
\]
Polarization of this equality gives the corresponding identity
for real inner products. Orthogonally to \(A'\) in \(p_\sigma Z\)
choose a vector of length \((30a^2+1)K\mu\), and take its inverse
image \(t\) in this pair summand. The choice makes \(t\) orthogonal
to every frequency of length less than \(\mu\), since its
projection already lies in \(A'\).

Choose a nearest \(u_s\in X\cap Z\) to \(t\). Take a real basis
of \(Z\) from the generators of length at most \(\Omega\mu\).
Rounding coordinates in its Gram--Schmidt directions, starting
with the last, gives
\[
   \|u_s-t\|\le\tfrac12\sqrt d\,\Omega\mu\le K\mu .
\]
Indeed at each step the residual orthogonal coordinate is at
most half the length of its Gram--Schmidt vector. The last
inequality follows from \(d\ge f\) and
\(d\Omega^2/4=df^3\rho^2\le d^4\rho^2\).
Define \(R_s=\|p_{(A')^\perp}p_\sigma u_s\|\) and \(e_s\)
to be that residual divided by \(R_s\). Then
\[
 R_s\ge30a^2K\mu,\qquad
 \|p_{A'}p_\sigma u_s\|\le K\mu,\qquad
 \|u_s\|\le\bigl(\sqrt2(30a^2+1)+1\bigr)K\mu\le45a^2K\mu .
\]
The first two estimates give \eqref{eq:light-triangular-projection};
they also prove that the induction increases the dimension.

For \(\|y\|<\mu\) compare \(u_s\) with \(u_s\pm y\in X\cap Z\).
Their distances from \(t\) and \((t,y)=0\) give
\(|2(u_s,y)|\le\|y\|^2\). For \(\|y\|\ge\mu\) the same
quantity is at most \(90a^2K\|y\|^2\). Finally use
\(P=2f/(\rho^2\lambda^2)\) in the thin range. Since
\(d^4=16f^4m^4\le16f^5\le8f\rho^3\), we get
\[
 (45a^2K\mu)^2\le
 \frac{2025}{2f}\,a^4d^8\rho^4P
 \le8100m^4d^4\rho^7P .
\]
\end{proof}

\begin{auxlemma}[One column with a prescribed short-period frequency]
\label{lem:light-frequency-column}
For \(u=u_s\) and \(R=R_s\) of
Lemma~\ref{lem:light-triangular-frequencies}, let \(N=2\pi R\)
and \(1\le c\le D\). In the theta frame at the origin there is a
level-\(M\) section \(s\) such that for unit vectors
\(v_j\in V_\sigma\),
\begin{align}
 \left|d^cs(0)[v_1,\ldots,v_c]
      -(-2\pi i)^c\prod_j H(v_j,u)\right|
      &\le 7\,2^g(4.3)^d e^{-T_0}N^c,
       \label{eq:light-frequency-jet}\\
 \|s\|_2^2 &\le K_6\exp(2\pi\|u\|^2/M)N^{2c},
 \qquad K_6=6(2d)^{d/2}.\nonumber
\end{align}
\end{auxlemma}
\begin{proof}
Choose \(x\in L\) with \(U(x)=u\) and put \(a=x/M\).
Use the lift of Lemma~\ref{lem:light-thick-lift} in coordinates
\((z,w)\in(W+iW)\oplus V_B\). Start with a quotient section
\(\beta_0\) whose Taylor polynomial through degree \(c\) is
\(1\). Use on the quotient the inverse unitary translation by
\(w_a\), the \(V_B\)-coordinate of \(a\), and choose its
semicharacter to give the required line after that translation.
The proof of Lemma~\ref{lem:light-thick-interpolation}
gives the uniform bounds
\[
 \|\beta_0\|_2^2\le 2M^{-\dim_{\mathbb C}V_B},\qquad
 \sup_w |\beta_0(w)|e^{-\pi M\|w\|^2/2}\le2 .
\]

Let \(z_a\) be the complex coordinates of the projection of
\(a\) to \(W+iW\). For the real orthonormal basis of \(W\) used
there, \(\operatorname{Im}z_a=u/M\). Apply the unitary translation
by \(a\) on the full torus after the lift, and divide by
\[
 C=\exp\!\left(\frac{\pi M}{2}
                 \left(\sum_j z_{a,j}^2-\|z_a\|^2\right)\right),
 \qquad |C|=e^{-\pi\|u\|^2/M}.
\]

The condition \(Ma\in L\) gives a translation of the same level
line. Multiplying its Gaussian factor by the translated zero mode
and completing the square gives
\[
   C^{-1}\,(\text{translated zero mode})(v)
    =e^{-2\pi i H(v,u)}
       \exp\!\left(\frac{\pi M}{2}\sum_j z_j(v)^2\right)
       \beta_0(w(v)).
\]

In fact the coefficient of the linear term in \(z_j(v)\) is
\(\pi M(\bar z_{a,j}-z_{a,j})\); the quotient translation has
cancelled the linear term in \(w(v)\). On \(V_\sigma\) the
quadratic form \(\sum_j z_j(v)^2\) is zero. It is the restriction
of the complexified Euclidean form on \(W\), invariant under the
real CM unit-circle action with characters \(e^{it},e^{-it}\)
on the conjugate pair. A square on the \(\sigma\)-eigenspace
has character \(e^{2it}\). The Taylor polynomial of \(\beta_0\)
now gives the main derivative in \eqref{eq:light-frequency-jet}.
The full translation is an isometry, so the norm comparison for the
lift and \(N\ge1\) give
\[
 \|s\|_2^2\le2|C|^{-2}\|\beta_0\|_2^2
             \le4e^{2\pi\|u\|^2/M}\le K_6e^{2\pi\|u\|^2/M}N^{2c}.
\]

It remains to estimate the nonzero modes after translation.
Put \(L_c=2\pi\|u\|+32\sqrt{M(c+1)}\) and \(b=c/L_c\).
For \(\|v\|\le b\) the term indexed by \(y\in X\setminus\{0\}\)
in the translated sum, divided by \(C\), has modulus at most
twice the exponential of
\[
 \tfrac12\pi Mb^2+2\pi b(\|u\|+M\|y\|)
               -\pi(M-M_1)\|y\|^2,\qquad
 M_1=90m^2d^2\rho .
\]
Indeed its Gaussian includes the difference
\(\pi M(\|u/M\|^2-\|\operatorname{Im}z(v)-u/M-y\|^2)\).
Expand the square, use \(\|\operatorname{Im}z(v)\|\le b\),
and apply \eqref{eq:light-frequency-cost} to the cross term
in \(u,y\). The remaining Gaussian at \(v\) costs at most
\(\pi Mb^2/2\). The factor two is the quotient supremum bound.

Using \(M_1\le M/2\), \(Mb^2\le c/1024\),
\(2\pi b\|u\|\le c\) and \(2b\|y\|\le16b^2+\|y\|^2/16\),
we bound the exponent by
\[
 c+\frac{33\pi c}{2048}-\frac{7\pi}{16}M\|y\|^2
       \le 1.06c-1.37M\|y\|^2 .
\]
Distinct frequencies are separated by \(\sqrt B\); disjoint
half-radius balls bound the count through radius \(t\sqrt B\)
by \((3t)^d\) for \(t\ge1\). Integrating this count against
\(2\alpha t e^{-\alpha t^2}\,dt\), where \(\alpha=1.37 T_0\),
and using \(2\log t\le t^2-1\) and \(T_0\ge d\), gives
\[
 \sum_{y\ne0}e^{-1.37M\|y\|^2}
   \le 3^d e^{-\alpha}\frac{\alpha}{\alpha-d/2}
   \le 2\,3^d e^{-1.37T_0}.
\]

For the mixed derivative use the polydisc
\(v=\sum_{j=1}^c t_jv_j\), \(|t_j|\le1/L_c\). The coefficient
of \(t_1\cdots t_c\) is the required derivative, so Cauchy's
formula costs exactly \(L_c^c\). Finally
\[
 L_c/N\le34\sqrt{M(c+1)},\qquad
 c\log\!\bigl(34e^{1.06}\sqrt{M(c+1)}\bigr)\le0.18T_0.
\]

The second inequality follows from \(c+1\le2c\) and
\(4c(7+\log(Mc))\le T_0\) in
\eqref{eq:joint-theta-log-cost}. The error is at most
\(4\,3^d e^{-1.19T_0}N^c\), which gives
\eqref{eq:light-frequency-jet}.
\end{proof}

\begin{auxlemma}[From estimates in one column to a determinant jet]
\label{lem:light-block-jet-product}
Let \(V\) have a Hermitian inner product \(H\), linear in the first
variable, and let \(\mathcal H\) be the Hilbert space of level-\(M\)
theta sections in one column with its Gaussian \(L^2\) norm. At the
origin use the common holomorphic theta frame. Partition some of the
\(n\) columns into disjoint ordered blocks \(B_j\) of size
\(1\le i_j\le r\). Fix orthonormal directions
\(e_{j1},\ldots,e_{ji_j}\) in the embedding space for each block.
For integers \(c_j\ge1\), write
\[
 \mathcal D=\prod_j
   \det\bigl(\partial_{e_{ja},k}\bigr)_{1\le a\le i_j,\,k\in B_j}^{c_j},
 \qquad D=\sum_j i_jc_j\ge1 .
\]

The second index of a derivative denotes its column. Let \(0<\omega\le1\)
and \(0\le\varepsilon\le\exp(-(D+1)(r\log r+3r+3))\).
For each \(k\in B_j\) suppose that there are
\(s_k\in\mathcal H\), \(\mu_k\ne0\), and \(y_k\in V\) such that
\begin{gather}
 \label{eq:light-single-column}
 \|y_k\|\le2,\qquad
 \omega\|s_k\|_2^2\le|\mu_k|^2,\qquad
 \left|\det(H(e_{ja},y_k))_{1\le a\le i_j,\,k\in B_j}\right|
       \ge\delta:=2^{-r/2}/2,\\
 \label{eq:light-single-column-derivative}
 \left|d^{c_j}s_k(0)[v_1,\ldots,v_{c_j}]
       -\mu_k\prod_l H(v_l,y_k)\right|\le\varepsilon|\mu_k|
\end{gather}
for unit directions in the indicated embedding space. In each unused
column suppose that \(s_k\ne0\) satisfies
\(\omega\|s_k\|_2^2/4\le|s_k(0)|^2\).
Then a unit vector \(F\in\mathcal H^{\otimes n}\) satisfies
\[
 |\mathcal D F(0)|^2\ge (\omega/4)^n 2^{-(r+2)D}.
\]
\end{auxlemma}
\begin{proof}
Take the product \(F_0=\bigotimes_k s_k\). In the determinant expansion
of \(\mathcal D F_0\), each term is a product of the derivatives of the
individual sections. The number of signed terms, before collecting
like terms, is
\[
 N=\prod_j(i_j!)^{c_j}\le r^D .
\]

Divide by \(s_k(0)\) in unused columns and by \(\mu_k\) in used ones.
The main term \(P_0\) and its lower bound are
\[
 P_0=\prod_j\det(H(e_{ja},y_k))^{c_j},\qquad
 |P_0|\ge\delta^{\sum_jc_j}\ge\delta^D .
\]
The bound on \(y_k\) shows that the normalized main factor in a column
of block \(j\) has modulus at most \(2^{c_j}\). Replacing it by the
actual jet changes it by at most \(\varepsilon\).
Let \(m_1=\sum_j i_j\le D\). In one signed product the error is at most
\[
 2^D\big((1+\varepsilon)^{m_1}-1\big)
       \le 2^D m_1\varepsilon(1+\varepsilon)^{m_1}
       \le D4^D\varepsilon .
\]

Multiplying by \(N\) still costs at most half the main term. Indeed
\[
 2D(4r)^D\delta^{-D}\varepsilon\le1.
\]

To check the last inequality, take logarithms and use
\(\log(2D)\le2D\), \(\log4<2\), and
\((r/2+1)\log2\le r/2+1\) in the assumed exponential bound
(or combine \(\log(2D)\le D\) and \(\log 4=2\log 2\)).
Thus
\[
 |\mathcal D F_0(0)|^2\ge
 \tfrac14\delta^{2D}
   \prod_{k\ {\rm used}}|\mu_k|^2
   \prod_{k\ {\rm unused}}|s_k(0)|^2 .
\]

There are \(m_1\ge1\) used columns, so the normalization gives
\[
 \tfrac14\omega^{m_1}(\omega/4)^{n-m_1}
       =4^{m_1-1}(\omega/4)^n\ge(\omega/4)^n .
\] Divide by
\(\|F_0\|_2^2=\prod_k\|s_k\|_2^2\) to obtain the claim.
\end{proof}

The norm on the tensor power also explains why this test bounds the
sum over the original marked translates. Let \(s\) generate the
principal theta line and let \(\Theta_\nu\) be the finite translates
of the pullback of \(s^{\boxtimes n}\) under the half-level isogeny.
The kernel acts unitarily. Its distinct characters give orthogonal
lines, and the translates run through them, by the character
decomposition in Lemma~\ref{lem:ideal-lattice-jet}. Their squared
norms are all \(h_s=\|s\|_2^{2n}\), because the Gaussian metric on the
pullback is the pullback metric and Haar probability pushes forward
to Haar probability. Expanding the unit vector of the lemma gives
\[
 F=\sum_\nu a_\nu\Theta_\nu,\qquad
 h_s\sum_\nu|a_\nu|^2=1,\qquad
 \sum_\nu|\mathcal D\Theta_\nu(0)|^2
       \ge h_s|\mathcal D F(0)|^2 .
\]
The last inequality is Cauchy--Schwarz. It is important here that the
test gives the size of a derivative of a \emph{unit} section.

To apply these estimates at the selected frame, set
\(d=2fm\), \(K_6=6(2d)^{d/2}\), and write \(P\) for the cusp
penalty of Lemma~\ref{lem:eta-period-decay}. Define
\[
 \omega=K_6^{-1}\exp\!\left(-\frac{2\pi}{M}
                 8100m^4d^4\operatorname{rd}(E)^7 P\right).
\]

There are enough distinct columns at this weight. With block ranks
\(p+q=r\), using a separate group of columns for every basic minor
requires at most
\[
 f\left\{\tbinom{p+1}2+\tbinom{q+1}2\right\}
 \le fr(r+1)/2\le n_W .
\]
Indeed the formula of Lemma~\ref{lem:rank-window} gives
\(n_W\ge10rf^2(k-1)^3L^2/m\), while \(m\le k\), \(r<2k\) and
\(k\ge5\). Repetitions of a basic minor are its exponent \(c_j\) in one
group of columns.

At each embedding take the triangular frequency directions of
Lemma~\ref{lem:light-triangular-frequencies} in the span of the short
periods, and complete by an orthonormal basis of its perpendicular
complement. Their one-column sections come respectively from
Lemmas~\ref{lem:light-frequency-column} and~\ref{lem:light-thick-lift},
on the same lattice of Lemma~\ref{lem:light-isotropic-periods}.
Cauchy--Schwarz gives the row-sum bound for the \(t\) short
directions in Lemma~\ref{lem:light-block-directions}:
\[
 \sum_{j=1}^t |E_{ij}|\le\sqrt{t}\,\|E_{i\bullet}\|_2
       \le \frac{\sqrt t}{30t^2}=\frac1{30t^{3/2}}\qquad(t>0).
\]
The frequency vectors used
after normalization have norm at most \(3/2\), and the perpendicular
directions have norm one. This verifies the determinant and
direction conditions in \eqref{eq:light-single-column}. For a thin column take \(\mu=(-2\pi iR)^c\), and for a perpendicular column take \(\mu=c!\). The two norm comparisons then give
\[
 \frac{\|s\|_2^2}{|\mu|^2}\le
 \begin{cases}
 K_6\exp(2\pi\|u\|^2/M),&\text{thin},\\
 K_6,&\text{perpendicular},
 \end{cases}
 \qquad
 \frac{\|s\|_2^2}{|\mu|^2}\le\omega^{-1}
\]
by \eqref{eq:light-frequency-cost}. For unused columns the
section constructed with the constant jet gives the stated
factor \(1/4\). Finally, \eqref{eq:joint-theta-rank-cost} gives
\[
 \log\varepsilon=\log7+g\log2+d\log(4.3)-T_0
 \le -(D+1)(r\log r+3r+3).
\]

Thus the marked sum has the lower bound of
Lemma~\ref{lem:light-block-jet-product} with this particular
\(\omega\).

\begin{lemma}[The scaled light generators]
\label{lem:scaled-light-metric}
The chosen scaled light generators span the integral light lattice.
On the framed generic fibre they have no common zero. For the generators \(w_l\) at a nondefinite lift,
\[
 \sum_l\|w_l\|^2\ge K_1^{-1}e^{-2cP}.
\]
In addition,
the scalar section of weight \(i_\eta g_0^*u_0\) in
\eqref{eq:selected-eta-product} is nonzero on every component and has squared norm at most \(K_2e^{-2i_\eta P}\) for the same
\(P\). At the definite lift the first-minimum penalty is zero.
\end{lemma}

\begin{proof}
\proofstep{Nonzero values for the scaled light family.}
Here is the nonvanishing argument in this normalization.
A triangular change to an orthonormal flag frame multiplies
every value at the same weight by its nonzero character.
Indeed, for the common weight character \(\chi\) and frame change \(b\),
\[
  (w_l(yb))_l=\chi(b)(w_l(y))_l,\qquad \chi(b)\ne0 .
\]
On the orthonormal frame apply Lemma~\ref{lem:ideal-lattice-jet}
with \(M_0\ge4D_F+4\) to the differential operator of a nonzero harmonic polynomial
of degree \(D_F\) in the prescribed flag type. The harmonic
polynomials used here exist by the central and flag weight
conditions. The frame-coordinate inclusion preserves
its leading tensor. The test gives a nonzero reading on
a theta section on the indicated sublattice. The finite
theta translation identity writes that section as a linear
combination of the allowed half-level translates.

Writing this translation identity as \(\Theta=\sum_\nu a_\nu\Theta_\nu\), the jet functional \(D\) from the harmonic polynomial therefore satisfies
\[
  0\ne D\Theta=\sum_\nu a_\nu D\Theta_\nu\quad\Longrightarrow\quad D\Theta_\nu\ne0\text{ for some }\nu.
\]
 Applying
the polynomial to this identity gives a nonzero reading
of some permitted theta rule. A representative of the
same level marking changes its reading by a nonzero
theta factor. This rule lies in the module spanned by
the unscaled family, so its multiple by \(s_W\) is a
linear combination of the chosen light columns.
Indeed, for a linear combination \(s_W H=\sum_j a_j w_j\), evaluation gives
\[
  (s_WH)(y)\ne0\quad\Longrightarrow\quad (w_j(y))_j\ne0 .
\]
 Some
column is nonzero. A common algebraic zero on a generic
framed chart would give a complex zero after the good
embedding, by faithful flatness and the Nullstellensatz.
The complex uniformization on the framed fibre excludes
it. This proves the required no-common-zero property on
the generic fibre.
\label{cor:slice-light-2}

\proofstep{Comparing the normalized light metric.}
For the metric comparison let \(P\) be the cusp penalty,
zero above the chart threshold for the first minimum and
equal below it to
\(2f/(\operatorname{rd}(E)^2\lambda_1^2)\).
Put \(M=g_0g_1\), \(D=D_F\ge1\), \(n=n(w_L)\), and
\(d_*=2fm\), where \(m\) bounds the smaller signature rank.
Define
\[
 B_*=24(2d_*)^{d_*/2},\qquad
 a=\frac{8100\pi n m^4 d_*^4\operatorname{rd}(E)^7}{M},\qquad
 J=B_*^{-n}2^{-(r+2)D}\exp(-2aP).
\]

Let \(b_\nu\) be the unscaled theta values for the translated block
determinant polynomial in the orthonormal frame. Use Lemma~\ref{lem:light-block-jet-product} at the parameters of
Lemma~\ref{lem:joint-theta-parameters}. Its one-column norm allowance is
\[
 \omega=\bigl[6(2d_*)^{d_*/2}\bigr]^{-1} e^{-2aP/n},
 \qquad (\omega/4)^n 2^{-(r+2)D}=J .
\]

The period comparison here uses the two elements \(\pi_i\) of
Proposition~\ref{prop:two-type-primes}. Write \(\alpha=\pi_0\pi_1\)
with the halves chosen by this orientation. The ideal computation in
Proposition~\ref{prop:scaled-integral-light-reading} gives
\(\mathfrak m=(\bar\alpha)\) and \(\alpha\bar\alpha=M\).
If \(H_A\) is the Hermitian form of the original principal polarization,
the principal form on the ideal lattice is \(H'=H_A/M\). Thus
\[
 (\Lambda,H_A)\longrightarrow(\mathfrak m\Lambda,H'),\qquad
 v\longmapsto\bar\alpha v,\qquad
 H'(\bar\alpha v,\bar\alpha v)=H_A(v,v).
\]

This map identifies the two first minima, the cutoff \(h\) and the
penalty \(P\). It also identifies the CM theta signs because it
commutes with the half-trace element. A marked unit letter has
squared \(H_A\)-length \(2\xi_R\) by definition of the polarization
scale; on the ideal torus its squared \(H'\)-length is \(2\xi_R/M\).
The unitary factors at each embedding preserve orthonormality.
Homogeneity of the order-\(D\) differential operator therefore
gives squared scale \((2\xi_R/M)^D\). Haar probability on the
ideal torus pushes forward under its isogeny to Haar probability
on the original torus, so the original theta section of unit
mean square contributes the factor one. The squared sum satisfies
\[
                    \sum_\nu |b_\nu|^2\ge (2\xi_R/M)^D J .
\]

The list of light generators includes these indexed translates after
multiplication by \(s_W\). Combining the bound
\(|\sigma(s_W)|^2\ge M^D\) with the squared differential normalization
\((2\pi\xi_R)^{-D}\) gives
\[
 S:=\sum_l\|w_l\|_{\rm rec}^2
 \ \ge\ (2\pi\xi_R)^{-D}M^D(2\xi_R/M)^D J=\pi^{-D}J.
\]

No minimum over the individual translates occurs here. Since \(J>0\),
taking logarithms gives
\[
 -\tfrac12\log S\le
 \tfrac12\big[D\log\pi+n\log B_*+(r+2)D\log2\big]+aP .
\]

The numerical allowances in the two runs are
\[
 A_5\ge\tfrac12\big[\log\pi+(r+2)\log2+n\log B_*\big],
 \qquad c\ge a .
\]

As \(D\ge1\), \(B_*>1\) and \(P\ge0\), this is at most \(A_5D+cP\).

If \(\ell_L\) converts the metric used in the height
calculation to \(\|\cdot\|_{\rm rec}\), the first-minimum
lower bound for the sum in that metric has
\(K_1=e^{2A_5D_F}\ell_L^2\). The scalar auxiliary form
is chosen independently of the later cut rows. Squaring the cusp estimate for this \(P\) gives
\[
 \ell_\eta^2\|\eta\|^2\le e^{-2i_\eta P}e^{2i_\eta A_2},
 \qquad
 \|\eta\|^2\le K_2e^{-2i_\eta P},\quad
 K_2=e^{2i_\eta A_2}\ell_\eta^{-2}.
\]
The two component bounds on \(\eta\) are its nonvanishing
as a section on each component and the displayed decay.

Strict dominance of \(i_\eta g_0^*w_L-w_\eta\)
comes from the weight gaps in the page. The inequality
\(2cg_0^*\le1\) ensures the strict cusp margin
\(c i_\eta g_0^*<i_\eta\).
The definite case uses \(P=0\), with no scalar auxiliary
form. These comparisons reduce the geometric light-section
inputs to the half-level jet and translated block estimates,
the auxiliary form and the specified page inequalities.
\label{cor:slice-light-3}
\end{proof}


\paragraph{The slice for the chosen sections.}
Write \(K=K_0\). A component \(Z\) of the moduli space over \(K\)
has a field \(\Omega_Z\) obtained by adjoining the coordinates of the
two open flag cells to its function field. If the signatures at a lift
are \((p,q)\), their dimensions and that of the underlying unitary
moduli space give
\[
 s=f\left(pq+\tbinom p2+\tbinom q2\right).
\]
Thus \(\Omega_Z/K\) is finitely generated of transcendence degree
\(s\). Evaluation in the generic flag frame sends a section to an
element of this field.

The components are finite in number and
open and closed. Indeed the generic fibre is smooth of finite type
over a field, and each of its local rings is regular and a domain.
Multiplying by the idempotent of one component extends a section
on it by zero.
If \(e_Z\) denotes this idempotent, its restrictions are
\[
  e_Z|_Z=1,\qquad e_Z|_{Z'}=0\quad(Z'\ne Z).
\]
 This preserves the weight and cuspidality, since the
indicator is locally constant on the analytic uniformization.
Vanishing on every generic flag cell detects the zero section by
density on the reduced charts.

Use the generators and the common first-minimum penalty from
the preceding calculation. The selected light generators span the scaled light lattice
and have no common zero on the framed generic fibre. Before the last
lift the first-minimum estimates on a good complex uniformization give
a function \(P\) and constants \(K_1,K_2>0\) such that
\[
 e^{-2cP}\le K_1 S,\qquad \|\eta\|^2\le K_2 e^{-2i_\eta P},
 \qquad S=\sum_l\|w_l\|^2.
\]
At sufficiently small first minimum \(\lambda_1\) we have
\(P\ge c_0/\lambda_1^2\), with \(c_0>0\).
The scalar section \(\eta\) is nonzero on each component.
With \(a=i_\eta g_0^*\), the run has \(ac<i_\eta\) and
\(a w_L-w_\eta\) strictly dominant in the flag and domain
directions. At the definite last lift use \(\eta=1\), \(a=1\)
and strict dominance of \(w_L\). We explain how these inputs
supply the generic dimension needed for the section calculation.
\label{cor:selected-slice-1}

\begin{proposition}[The reduced generic intersection and evaluation on its branches]
\label{prop:selected-generic-cut}
The generic cut defined by the chosen light rows is finite and reduced.
It is computed from the product of the component fields by the
equal-weight ratios of the light sections. At every branch where
the reference section is nonzero, the independent zeroth row attains
the minimum of the light-column orders. If \(v\) is the branch order,
\(e\) its relative ramification and \(c_\zeta=v(g_0(y_\zeta))\), then
for weight \(w_F+Tw_L\)
\[
 v\big((F/(g_0L_0^T))|_\zeta\big)+c_\zeta
      =e\big(v_{y_\zeta}(F)-T\nu_\zeta\big),
 \qquad \nu_\zeta=v(L_0(y_\zeta))/e.
\]
\end{proposition}

\begin{proof}
\proofstep{Finding a generic slice with the chosen sections.}
At a point with \(\|\eta\|^2/S^a\ge d>0\), substitute the bounds for
\(S^{-1}\) and \(\|\eta\|^2\):
\[
 d\le \|\eta\|^2S^{-a}\le K_2 e^{-2i_\eta P}K_1^ae^{2acP},
 \qquad
 P\le\frac{\log(K_1^aK_2/d)}{2(i_\eta-ac)} .
\]
Here the strict margin \(ac<i_\eta\) permits division.
There the first minimum is bounded below. Reduction theory gives
compact representatives up to the arithmetic group and changes
of flag frame. The definite space already has compact representatives.
The invariant ratio therefore attains a positive maximum.
Writing \(\Phi=\|\eta\|^2/S^a\), the maximizing point \(x_0\) thus satisfies
\[
  \Phi(x_0)=\max_x\Phi(x)>0,\qquad S(x_0)>0 .
\]
 Choose a
nonvanishing column there.

In holomorphic frames the Levi form of
the logarithm of the invariant ratio is the curvature of
\(a w_L-w_\eta\) minus \(a\) times the pullback of the
Fubini--Study form under \(x\mapsto[w_l(x)]\). This follows by
writing the logarithm as the sum of \(\log|\eta|^2\), the weight
potential, and \(-a\log\sum_l|w_l|^2\). The first summand is
pluriharmonic near the maximum. The invariant metric on a strictly
dominant weight has positive Levi form on these tangent directions
(the coefficients on the flag roots are the adjacent weight
differences).

Write \(\Phi\) for the invariant ratio. Along a nonzero vector \(v\) killed by the projective differential, the pulled-back projective form vanishes. The remaining strictly positive term would give
\[
  {\cal L}\log\Phi(v,\bar v)>0,
\]
contradicting negative semidefiniteness at the maximum. The projective differential is therefore injective. Some \(s\) ratios
to the nonvanishing column then have an open set of complex values.
A polynomial relation over \(K\) on \(Z\), homogenized using this
column, would vanish on the corresponding analytic component.
It follows that the field of the light ratios has transcendence
degree \(s\). The compact set and the maximizing point here may
depend on the component and weight. Only the dimension conclusion
is used.

\paragraph{Constant-coefficient ratios.}

Before eliminating the slice variables, we need to know the field generated by the ratios of the chosen light sections on each component.

Fix a column \(w_b\) nonzero at the generic frame of \(Z\).
The light ratios and the component field have the same transcendence degree. Additivity therefore gives
\[
  \operatorname{trdeg}_{K(w_l/w_b)}\Omega_Z=s-s=0.
\]
Since \(\Omega_Z\) is finitely generated over the field of ratios, this algebraic extension is finite. Let \(u\in\Omega_Z\)
have a monic equation of degree \(n\) over the ratio field. Each
coefficient \(a_i\) for \(0\le i<n\) is a quotient
\(A_i(w)/B_i(w)\) of homogeneous polynomials of the same degree.
Put 
\[
  P(w)=w_b\prod_{i<n} B_i(w),
\]
 of degree \(k>0\).
Then \(z=uP(w)\) satisfies
\[
 z^n+\sum_{i<n} C_i(w)z^i=0,\qquad
 C_i=P^{n-i}A_i/B_i,\qquad \deg C_i=(n-i)k .
\]
All \(C_i\) are polynomials.

In each flag chart the weight
transition of \(z\) is the \(k\)-th power of that of the light
weight. The transformed equation is monic with regular coefficients.
The chart ring is smooth and integrally closed in its fraction field.
Thus the transformed rational function is regular on the chart.
These sections agree on overlaps by agreement at the generic point.
Extend by zero to the other components.

\proofstep{Returning the ratio to a cusp section.}
The resulting weight-\(k w_L\) section \(g\) is cuspidal. For an
explicit check the homogeneous root bound in a normalized frame
gives \(\|g\|\le C(\max_l\|w_l\|)^k\): divide the equation by
the corresponding power of this maximum and apply the bound for
roots of a monic polynomial with bounded coefficients. Where all
coefficients vanish the root is zero. Multiplying the monic relation by \(g\) gives
\[
               g^{n+1}=-\sum_{i=0}^{n-1} C_i(w)g^{i+1}
\]
on all components, including the components where \(g=0\).
Each coefficient \(C_i(w)\) has positive degree \((n-i)k\)
in cusp forms, so every summand on the right is cuspidal.

On a rational boundary face, the allowable Fourier indices
are in the positive semidefinite cone. A sum of indices lies
in the face only if each index lies in it. Let \(\operatorname{CT}_{\mathfrak f}\) denote restriction to
the indices of this face \(\mathfrak f\). A product can contribute
there only through factors whose indices all lie on the face.
Consequently
\[
 \bigl(\operatorname{CT}_{\mathfrak f}g\bigr)^{n+1}
   =\operatorname{CT}_{\mathfrak f}(g^{n+1})
   =-\sum_{i=0}^{n-1}
     \operatorname{CT}_{\mathfrak f}(C_i(w))\,
           \bigl(\operatorname{CT}_{\mathfrak f}g\bigr)^{i+1}
   =0.
\]
Every \(C_i(w)\) is cuspidal. The boundary reading is over a reduced
complex chart, so \(\operatorname{CT}_{\mathfrak f}g=0\). We have obtained
\(u=g/P(w)\) at the generic frame. Both numerator and denominator
may be padded to any larger degree using \(w_b\); on a finite set
of components take the largest degree and extend by zero before
summing. No degree bound from this auxiliary ratio calculation
enters the index. \label{cor:selected-slice-2}

\paragraph{The algebra of the generic intersection.}

With these ratios fixed, the slice equations eliminate one coefficient in every row. This gives the algebra on which branch orders and trace discriminants will be computed.

Let \(F=K(X)\), with all \((s+1)N\) coefficients independent.
On a component put \(x_l=w_l/w_b\) and eliminate \(X_{ib}\) for
\(1\le i\le s\) using 
\[
  X_{ib}=-\sum_{l\ne b}X_{il}x_l.
\]
Denote the other variables by \(X'\), including every variable of
row zero. The substitution embeds \(F\) in \(L_Z=\Omega_Z(X')\).
Indeed the differentials of \(s\) of the \(x_l\) are independent
in characteristic zero.

Let \(l_1,\dots,l_s\) index such a set. Modulo \(dX'\) the
determinant test is
\[
 \left[\prod_{i=1}^s X_{i l_i}\right]
      \bigwedge_{i=1}^s\!\left(-\sum_{l\ne b}X_{il}\,dx_l\right)
       =(-1)^s dx_{l_1}\wedge\cdots\wedge dx_{l_s}\ne0 .
\]
The variables of different cut rows are independent. Hence the
full substituted coefficient family is algebraically independent.
The field \(L_Z\) is finitely generated, and its transcendence
degree over \(K\) is 
\[
  s+((s+1)N-s)=(s+1)N.
\]
 Thus \(L_Z/F\)
is finite, and separability follows from characteristic zero.
More precisely in \(\Omega_Z[X]\) the kernel of substitution is
the ideal of the \(s\) row equations, by division by their pivot
variables. After localizing by \(K[X]\setminus\{0\}\) the quotient
is a domain algebraic over \(F\), hence a field with fraction field
\(L_Z\). The independent zeroth row has nonzero value there.
After the coefficient extension to \(F'=K_{0,\lambda}(X)\) the
component algebra is 
\[
  F'\otimes_F L_Z,
\]
 finite and reduced.

\proofstep{Passing from fields to the actual cut.}
We verify why this calculation in function fields applies to the
weight-zero chart of the original forms. Evaluation embeds the chart
before cutting in the product of the component chart rings
\((F'\otimes_K\Omega_Z)[1/L_0]\); injectivity follows by putting
the fractions over a common power of \(L_0\), using independence
of scalar coefficients and detection of a homogeneous numerator at
the generic frames. For each component its indicator is in the chart:
replace each \(w_l\) in the zeroth row by its extension by zero from
that component and divide the result by \(L_0\).
If \(w_{l,Z}\) is the extension by zero of \(w_l|_Z\), the component indicator is
\[
  e_Z=\frac{\sum_lX_{0l}w_{l,Z}}{L_0},\qquad e_Z|_{Z'}=\begin{cases}1&Z'=Z,\\0&Z'\ne Z.\end{cases}
\]
 These orthogonal
idempotents split the chart.

On one summand let \(Q\) be any prime containing the cut equations,
and read the light columns as \(w_l/L_0\) in its residue field. Some
\(w_b/L_0\) is nonzero there by the zeroth-row identity. Invert it.
The algebra of constant homogeneous numerator ratios with pivot
denominator \(w_b^k\) lies in \(\Omega_Z\) and specializes at \(Q\).
Let \(R\subset\Omega_Z\) be that algebra of ratios and
\(\bar R\) its image at \(Q\). The original independent
coefficients are recovered from \(\bar R\) and the free
coefficients by the pivot equations. Therefore
\[
 (s+1)N
 \le \operatorname{trdeg}_K\operatorname{Frac}(\bar R)
                            +((s+1)N-s),
\]
while \(\operatorname{trdeg}_K\Omega_Z=s\). The two bounds
force the residue transcendence degree to be \(s\).
The specialization is then injective, as follows.

To see this, choose \(s\) ratios with
algebraically independent residues (the same dimension count applies
to the algebra generated just by the light ratios). Invert all
nonzero polynomials in them. The localized domain is algebraic over
the resulting rational function field and hence a field.
Indeed a nonzero element \(a\) algebraic over the rational function field has a monic minimal polynomial with \(c_0\ne0\), so
\[
  a^{-1}=-c_0^{-1}\big(a^{d-1}+c_{d-1}a^{d-2}+\cdots+c_1\big).
\]
 Its map to
the residue field has zero kernel. Thus every nonzero homogeneous
denominator of the constant-coefficient ratio construction is
invertible in the local ring at \(Q\), after dividing by a power of
\(L_0\).

\proofstep{Checking that the two localizations agree.}
Consequently the component chart, localized near \(Q\), already
contains all functions of \(\Omega_Z\) as ratios in weight zero.
It also recovers the inverse of the generic zeroth row in the
chosen generic frame: write \(1/w_b\) as a component function
and multiply by \(w_b/L_0\).
In the chosen frame the equality used here is
\[
  (1/w_b)(w_b/L_0)=1/L_0 .
\]
 Conversely evaluation embeds this
localization in the corresponding localization of the component
chart ring. More explicitly, localize at the elements of the
weight-zero chart which are nonzero at every prime of its slice.
The same denominator argument gives \(\Omega_Z\) there, using a
fixed generically nonzero \(w_b\), which stays nonzero at every
such prime by the specialization just proved. The constants, the
coefficients and \(1/L_0\) generate the component chart ring.
Upon cutting, all the inverted denominators were already units.

In the other direction these denominators stay units on the cut
of the component chart ring, since a prime there contracts to a
prime containing the cuts. This identifies the actual slice with
\[
                   \prod_Z (F'\otimes_F L_Z).
\]
In this description \(L_0\) in the component chart ring denotes
its generic-frame reading. The independent frame scalings disappear
in its weight-zero fractions.

\paragraph{Integral frames and the order at a branch.}

A local order is an order of a section in a chosen frame. The cusp form and its complementary section furnish such frames for the branches of the slice.

For the weight-\(w_F\) sheaf the run uses a cusp form \(g_0\)
nonzero on every component and a form \(t_2\) of complementary
weight \((s+4)w_L-w_F\), also generically nonzero on each.
For \(g_0\), nonzero values of the theta sections of weight \(w_F\) provide
a proper vanishing subspace on each of the finitely many analytic
components. Their union cannot cover a vector space over the
infinite coefficient field. The level test here uses the positive
last weights and 
\[
  3\operatorname{DF}(w_F)+2\le M
\]
 where \(M\)
is the full level. Multiplication by \(g_0\) is injective by generic
evaluation. The degree-zero element
\(b=g_0t_2/L_0^{s+4}\) is invertible on the slice just identified.
Indeed it is a unit on each incidence field and stays a unit after
extension of scalars.

Take a chart element \(b'\) with 
\[
  b'b=1\pmod J.
\]
 For
\(a\ge1\), an inverse on the thickened chart is
\[
    b^{-1}\equiv b'\sum_{k=0}^{a-1}(1-b'b)^k\pmod{J^a},
 \qquad
 \frac{h}{g_0}
  =b^{-1}\frac{h t_2}{L_0^{s+4}}\pmod{J^a}
       \quad(h\text{ of weight }w_F).
\]
Indeed \((1-b'b)^a\in J^a\), and the numerator in the second
formula has weight zero after division by the indicated power.
This gives the asserted trivialization by \(g_0\) on the cut
and its positive-order thickenings.

\proofstep{The branch order in the reference frame.}
At a branch \(\zeta\) take the completion of its normalization
DVR. Write \(U\) for the fraction field, \(v\) for its order
and \(e=v(\pi)\). The embeddings 
\[
  \Omega_Z\to L_Z\to U
\]
and \(F'\to U\) evaluate the original rows. Over \(U\) pull
the generic framed point on \(Z\) through the first embedding.
The integral model used for this operation is the proper
compactification with its semiabelian family and invertible
pushforward of the theta line. Its good-level condition at a
residue prime is that the prime-to-residue-characteristic part
of the level is at least three. The two different odd level
primes ensure it even in their own residue characteristics.
The valuation ring of the coefficient model maps into the
branch DVR since the branch order restricts to \(e v_\lambda\).

In particular nonnegative coefficient orders stay nonnegative at the branch:
\[
  v_\lambda(a)\ge0\quad\Longrightarrow\quad v(a)=e\,v_\lambda(a)\ge0.
\]
The valuative criterion extends the underlying point to that DVR.
In the pulled-back Lie lattices the intersections with the
successive generic flag subspaces are saturated. Their successive
quotients are free over the DVR, so choose compatible integral
bases. Choose also a generator of the pulled-back invertible
theta module. These choices reframe the original field point
by an upper triangular change in each flag and a scalar in the
theta line. If \(A\) denotes the block triangular and theta change,
a section \(g\) of weight \(u\) has reading

\[
  g(y_\zeta)=\chi_u(A)g|_\zeta.
\]
 The multiplier is the
same for every section of that weight. A second integral
basis for the same pulled-back flag and theta line changes
it by a unit, since its diagonal entries and theta scalar
are units. Hence for equal-weight \(g,h\) with nonzero
denominator,
\[
   v\big(g(y_\zeta)/h(y_\zeta)\big)
         =v\big((g/h)|_\zeta\big).
\]
This places the branch calculation in the extended set
of integral tests. It gives nonzero
evaluations of \(g_0\) and, before the final lift, \(\eta\).

\proofstep{The minimum from the independent zeroth row.}
The independent zeroth row also computes the minimum on the
light lattice. Choose a column \(b\) of smallest finite order
on \(Z\) at this branch. In the pivot elimination in the next
paragraph the free variables have Gauss valuation over the valued
component field, including all the variables of row zero.
It follows that
\[
 v\left(\sum_l X_{0l}w_l/w_b\right)=
 \min_l v(w_l/w_b)=0 .
\]

For clarity, this uses only the residue transcendence-degree
calculation below, which needs no minimum assertion about \(L_0\).
The light columns span the scaled light lattice over the integers.
The ultrametric inequality and attainment at \(b\) identify its
depth at \(y_\zeta\) with \(\nu_\zeta=v(L_0(y_\zeta))/e\).
Similarly \(c_\zeta=v(g_0(y_\zeta))\). Writing
\(v_y=v/e\), we have for weight \(w_F+Tw_L\)
\[
 v\big((F/(g_0 L_0^T))|_\zeta\big)+c_\zeta
                   = e\big(v_{y_\zeta}(F)-T\nu_\zeta\big).
\]
All entries on the left are integer orders or infinity. In particular
each integral source section has nonnegative unshifted point valuation
by the definition of the cusp lattice on integral frames. These are
the same points for the branch length, the depth test and the lifting
through \(J\). \label{cor:selected-slice-3}
\end{proof}
